\documentclass[../../../include/open-logic-section]{subfiles}

\begin{document}

\olfileid[id]{sth}{spine}{recursion}
\olsection{Teorema(-Teorema) Rekursi Transfinit}

Namun, hal pertama yang harus kita lakukan ialah memastikan bahwa
\olref[valpha]{defValphas} merupakan definisi yang berhasil. Secara lebih
umum, kita perlu membuktikan bahwa setiap upaya memberikan suatu definisi
transfinit melalui rekursi (transfinit) akan berhasil. Itulah tujuan bagian
ini.

\emph{Peringatan: materi ini rumit}. Meskipun demikian, pesan utamanya cukup
sederhana: Induksi Transfinit bersama Penggantian menjamin keabsahan
(beberapa versi) rekursi transfinit.\footnote{Pengingat: semua formula dan
suku dapat mempunyai parameter (kecuali dinyatakan sebaliknya secara
eksplisit).}

\begin{defn}
	Misalkan $\tau(x)$ suatu suku; misalkan $f$ suatu fungsi; misalkan $\alpha$
	suatu ordinal. Kita mengatakan bahwa $f$ merupakan suatu
	\emph{aproksimasi-$\alpha$} bagi $\tau$ jika dan hanya jika berlaku baik
	$\dom{f} = \alpha$ maupun $(\forall \beta \in \alpha)f(\beta) =
	\tau(\funrestrictionto{f}{\beta})$.
\end{defn}

\begin{lem}[Rekursi Terbatas]\ollabel{transrecursionfun}
Untuk sembarang suku $\tau(x)$ dan sembarang ordinal $\alpha$, terdapat suatu
aproksimasi-$\alpha$ yang unik bagi $\tau$.
\end{lem}
\begin{proof}
Kita akan menunjukkan bahwa, untuk setiap $\gamma \leq \alpha$, terdapat
suatu aproksimasi-$\gamma$ yang unik.

Pertama-tama kita menetapkan keunikan. Misalkan $g$ dan $h$ (berturut-turut)
merupakan aproksimasi-$\gamma$ dan aproksimasi-$\delta$. Induksi transfinit
pada argumen keduanya menunjukkan bahwa $g(\beta) = h(\beta)$ untuk setiap
$\beta \in \dom{g} \cap \dom{h} = \gamma \cap \delta = \min(\gamma,
\delta)$. Jadi, aproksimasi kita unik (jika ada), dan semuanya bersepakat
pada semua nilai.

Untuk menetapkan keberadaan, kini kita menggunakan induksi transfinit
sederhana (\olref[ordinals][opps]{simpletransrecursion}) pada ordinal
$\delta \leq \alpha$.

Fungsi kosong secara trivial merupakan suatu aproksimasi-$\emptyset$.

Jika $g$ merupakan suatu aproksimasi-$\gamma$, maka $g \cup
\{\tuple{\gamma, \tau(g)}\}$ merupakan suatu
aproksimasi-$\ordsucc{\gamma}$.

Jika $\gamma$ merupakan suatu ordinal limit dan $g_\delta$ merupakan suatu
aproksimasi-$\delta$ untuk semua $\delta < \gamma$, misalkan $g =
\bigcup_{\delta \in \gamma} g_\delta$. Ini merupakan suatu fungsi karena
berbagai $g_\delta$ kita bersepakat pada semua nilai. Jika $\delta \in
\gamma$, maka $g(\delta) = g_{\ordsucc{\delta}}(\delta) =
\tau(\funrestrictionto{g_{\ordsucc{\delta}}}{\delta}) =
\tau(\funrestrictionto{g}{\delta})$.

Ini menyelesaikan pembuktian melalui induksi transfinit.
\end{proof}

Jika kita membolehkan diri mendefinisikan suatu \emph{suku} alih-alih suatu
fungsi, maka kita dapat menghapus batas $\alpha$ dari hasil sebelumnya.
Dalam pernyataan dan pembuktian hasil berikut, jika $\sigma$ merupakan suatu
suku, kita misalkan $\funrestrictionto{\sigma}{\alpha} =
\Setabs{\tuple{\beta, \sigma(\beta)}}{\beta \in \alpha}$.

\begin{thm}[Rekursi Umum]\ollabel{transrecursionschema}
Untuk sembarang suku $\tau(x)$, kita dapat secara eksplisit mendefinisikan
suatu suku $\sigma(x)$ sedemikian sehingga $\sigma(\alpha) =
\tau(\funrestrictionto{\sigma}{\alpha})$ untuk sembarang ordinal~$\alpha$.
\end{thm}

\begin{proof}
Untuk setiap $\alpha$, menurut \olref{transrecursionfun} terdapat suatu
aproksimasi-$\alpha$ yang unik, $f_\alpha$, bagi~$\tau$. Definisikan
$\sigma(\alpha)$ sebagai $f_{\ordsucc{\alpha}}(\alpha)$. Sekarang:
	\begin{align*}
		\sigma(\alpha) &=
		f_{\ordsucc{\alpha}}(\alpha) \\&=
		\tau(\funrestrictionto{f_{\ordsucc{\alpha}}}{\alpha}) \\&=
		\tau(\Setabs{\tuple{\beta, f_{\ordsucc{\alpha}}(\beta)}}{\beta \in \alpha}) \\&=
		\tau(\Setabs{\tuple{\beta, f_{\ordsucc{\beta}}(\beta)}}{\beta \in \alpha})\\&=
		\tau(\funrestrictionto{\sigma}{\alpha})
	\end{align*}
dengan memperhatikan bahwa $f_{\ordsucc{\beta}}(\beta) =
f_{\ordsucc{\alpha}}(\beta)$ untuk semua $\beta < \alpha$, sebagaimana
dalam \olref{transrecursionfun}.
\end{proof}
\noindent
Perhatikan bahwa \olref{transrecursionschema} merupakan suatu \emph{skema}.
Yang penting, kita tidak dapat mengharapkan $\sigma$ mendefinisikan suatu
fungsi, yakni suatu jenis \emph{himpunan} tertentu, sebab dalam hal itu
$\dom{\sigma}$ akan menjadi himpunan semua ordinal, bertentangan dengan
Paradoks Burali-Forti (\olref[ordinals][basic]{buraliforti}).

Namun, masih perlu ditunjukkan bahwa \olref{transrecursionschema}
mengabsahkan definisi kita atas $V_\alpha$. Hal ini mungkin tidak langsung
jelas; tetapi akan tampak dengan suatu versi terakhir yang sederhana dari
rekursi transfinit.

\begin{thm}[Rekursi Sederhana]\ollabel{simplerecursionschema}
Untuk sembarang suku $\tau(x)$ dan $\theta(x)$ serta sembarang himpunan $A$,
kita dapat secara eksplisit mendefinisikan suatu suku $\sigma(x)$ sedemikian
sehingga:
\begin{align*}
	\sigma(\emptyset) &= A\\
	\sigma(\ordsucc{\alpha}) &= \tau(\sigma(\alpha)) &&
		\text{untuk sembarang ordinal }\alpha\\
	\sigma(\alpha) &= \theta(\ran{\funrestrictionto{\sigma}{\alpha}})&&
	\text{ketika }\alpha\text{ merupakan ordinal limit}
\end{align*}
\end{thm}

\begin{proof}
Kita mulai dengan mendefinisikan suatu suku, $\xi(x)$, sebagai berikut:
%t $\xi(x) = A$ if $x$ is not a function whose domain is an ordinal;
%otherwise let:
\[
	\xi(x) =
	\begin{cases}
			A &\text{jika $x$ bukan fungsi yang domainnya}\\
			&\hspace{1em}\text{ordinal tak kosong; jika tidak:}\\
			\tau(x(\alpha)) & \text{jika $\dom{x} = \ordsucc{\alpha}$}\\
			\theta(\ran{x}) & \text{jika $\dom{x}$ merupakan ordinal limit}
	\end{cases}
\]
Menurut \olref{transrecursionschema}, terdapat suatu suku $\sigma(x)$
sedemikian sehingga $\sigma(\alpha) =
\xi(\funrestrictionto{\sigma}{\alpha})$ untuk setiap ordinal $\alpha$;
selain itu, $\funrestrictionto{\sigma}{\alpha}$ merupakan suatu fungsi dengan
domain $\alpha$. Kita menunjukkan bahwa $\sigma$ mempunyai sifat-sifat yang
diperlukan, dengan induksi transfinit sederhana
(\olref[ordinals][opps]{simpletransrecursion}).

Pertama, $\sigma(\emptyset) = \xi(\emptyset) = A$.

Berikutnya, $\sigma(\ordsucc{\alpha}) =
\xi(\funrestrictionto{\sigma}{\ordsucc{\alpha}}) =
\tau(\funrestrictionto{\sigma}{\ordsucc{\alpha}}(\alpha)) =
\tau(\sigma(\alpha))$.

Terakhir, $\sigma(\alpha) =
\xi(\funrestrictionto{\sigma}{\alpha}) =
\theta(\ran{\funrestrictionto{\sigma}{\alpha}})$, jika $\alpha$ merupakan
suatu limit.
\end{proof}
\noindent
Sekarang, untuk mengabsahkan \olref[valpha]{defValphas}, cukup ambil $A =
\emptyset$ dan $\tau(x) = \Pow{x}$ serta $\theta(x) = \bigcup x$. Akhirnya,
ini mengabsahkan definisi $V_\alpha$!{}


\end{document}
